\section{A complete six-rule presentation}\label{sec:system}
\subsection{Syntax and computation}
There is one numerical sort with raw terms
\[
 t::=x\mid0\mid S(t)\mid I(n,a;u.h).
\]
Only $u$ in $h$ is bound. Updates are fixed finite numerical templates
with passive parameters; there are no function variables or universal
evaluators for primitive-recursive descriptions.

Represent terms in full normal form for
\[
 I(0,a;u.h)\longrightarrow a,\qquad
 I(Sn,a;u.h)\longrightarrow h[I(n,a;u.h)/u].
 \tag{Comp}\label{eq:comp}
\]
Every substitution is capture avoiding and followed by normalization.
For $C(v,a,b):=I(v,a;u.b)$ with unused $u$,
\[
 C(0,a,b)\equiv a,\quad C(Sv,a,b)\equiv b,\quad
 I(1,a;u.h)\equiv h[a/u].
\]
Definitional equality here uses only the two computation equations,
not arbitrary proved arithmetic equalities.

Normalization is counted as part of the trusted specification.
The repository supplies totality, canonical forms and
\[
 \nf(\nf(t)[\sigma])=\nf(t[\sigma]).                    \tag{NCoh}
\]
Its normalization-by-evaluation proof interprets a structurally smaller
body as a total operation on normal terms, traverses the finite
successor prefix of the normalized counter, and retains an iterator at
a neutral counter. Environment laws establish substitution coherence.
The reference checker uses equivalent direct reductions with capture
avoidance and memoization \cite{normalization,current}.

Formulas and their ordinary readings are
\[
 A::=(s=t)\mid\Q_yA\mid\D_{x,y}A,
\]
\[
 \Q_yB:=B(0)\to\forall y\,B(y),\qquad
 \D_{x,y}B:=\exists x\bigl(B(x,0)\to\forall y\,B(x,y)\bigr).
 \tag{Read}\label{eq:read}
\]
D binds distinct $x,y$. Contexts are finite reusable lists of formulas.
The native contradiction is $0=S0$. The readings do not silently insert
ordinary equality logic.

\subsection{All primitive rules}
All displayed expressions are normalized. The six families are
\[
 \frac{A\in\Gamma}{\Gamma\vdash A}\ {\rm Hyp}
 \qquad
 \frac{\ }{\Gamma\vdash t=t}\ {\rm Ref},
\]
\[
 \frac{\Gamma\vdash\Q_yB(y)\qquad\Gamma\vdash B(0)}
      {\Gamma\vdash B(t)}\ {\rm QE},
\]
\[
 \frac{\Gamma\vdash\D_{x,y}B(x,y)\qquad
       \Gamma,\Q_yB(b,y)\vdash A}
      {\Gamma\vdash A}\ {\rm DE},
\]
where $b$ is fresh for $\Gamma,A$ and the passive parameters of $B$;
\[
 \frac{\Gamma,B(t,0)\vdash B(t,Sz)}
      {\Gamma\vdash\D_{x,y}B(x,y)}\ {\rm DI_c},
\]
where $z$ is fresh for $\Gamma,t$ and the passive parameters of $B$; and
\[
 \boxed{\quad
 \frac{\Gamma\vdash f(0)=a\qquad
       \Gamma\vdash f(Sz)=h(f(z))}
      {\Gamma\vdash I(n,a;u.h(u))=f(n)}\ \U .
 \quad}                                                \tag{U}\label{eq:U}
\]
In U, $f(z),h(u)$ are fixed finite numerical templates. The eigenparameter
$z$ is fresh for $\Gamma,a,n,h$ and all passive parameters of $f$.
Thus $h$ has no active-counter argument. The recurrence premise has no
induction hypothesis. DI$_c$ always uses cases, even at equation matrices.

There are no other logical axioms, computation axioms, transport rules
or induction rules. Computation is exactly \eqref{eq:comp} in the syntax.

\begin{lemma}[Structural admissibility]
Weakening, renaming, ordinary-assumption cut and numerical substitution
followed by normalization are admissible.
\end{lemma}
\begin{proof}
Induct on finite derivations, freshening eigenparameters first.
Normalization coherence identifies substituted premises and conclusions.
Neither DI$_c$ nor U has a mode that can change under substitution.
Replace each use of a cut assumption by its proof. These are
metatheoretic inductions on finite syntax, not numerical induction.
\end{proof}

\section{Equality derived before induction}\label{sec:equality}
\subsection{Numerical cases without transport}
\begin{lemma}
The case logic, without U, derives
\[
 \frac{\Gamma\vdash A(0)\qquad\Gamma\vdash A(Sz)}
      {\Gamma\vdash A(n)}
 \quad\text{and}\quad
 \frac{\Gamma,A(0)\vdash A(Sz)}
      {\Gamma\vdash\Q_yA(y)},                          \tag{Cases}\label{eq:cases}
\]
with fresh $z$.
\end{lemma}
\begin{proof}
DI$_c$ with an unused coordinate $w$ introduces $\D_{w,y}A(y)$.
For the first rule weaken the successor proof by $A(0)$.
DE supplies $\Q_yA(y)$, and QE with the base gives $A(n)$,
independent of the opened witness. For the second rule DE exports the
opened Q assumption itself by Hyp.
\end{proof}

\subsection{Symmetry}
Given $s=t$, choose
\[
 f(z)=s,\qquad a=t,\qquad h(u)=u,\qquad n=0
\]
in U. Its recurrence is Ref. Its conclusion normalizes to $t=s$.

\subsection{Congruence in any numerical template}
For a template $H(u)$ and given $s=t$, choose
\[
 f(z)=I(z,s;u.H(u)),\qquad a=t,\qquad h=H,\qquad n=1.
\]
The base is $s=t$, while both sides of the recurrence normalize to
$H(I(z,s;H))$. Ref proves it. The conclusion is $H(t)=H(s)$;
symmetry gives ordinary congruence. Templates may duplicate or omit
the argument, or place it beneath additional iterator binders.

\subsection{Joining and transitivity}
\begin{lemma}[Common-source joining]
From $s=t$ and $s=u$ in one retained context, derive $t=u$ without
numerical induction.
\end{lemma}
\begin{proof}
Let $f(z)=C(z,s,u)$ and use the identity update. Its recurrence is
\[
 u=C(z,s,u).
\]
At zero this is $u=s$, by symmetry of $s=u$; at a successor it is
$u=u$, by Ref. Numerical cases prove the recurrence.
U with $a=t,n=1$ has base $s=t$ and conclusion $t=u$.
\end{proof}
From $s=t,t=u$, first derive $t=s$ and join at $t$ to get $s=u$.
Thus transitivity is derived. From $s=t,H(s)=r$, congruence supplies
$H(s)=H(t)$ and joining gives $r=H(t)$, the repository's one-term
crossed transport rule.

\subsection{Full formula transport}
\begin{theorem}[Leibniz transport]
For every native formula template $F(v)$,
\[
 \Gamma\vdash s=t,\quad\Gamma\vdash F(s)
       \quad\Longrightarrow\quad\Gamma\vdash F(t).
\]
\end{theorem}
\begin{proof}
Induct on $F$. At an equation use congruence on both sides, symmetry and
transitivity.

For $F(v)=\Q_yB(v,y)$, temporarily assume $B(t,0)$.
Reverse matrix transport gives $B(s,0)$; QE at $Sz$ gives $B(s,Sz)$;
forward matrix transport gives $B(t,Sz)$. Conditional Q introduction
\eqref{eq:cases} finishes.

For $F(v)=\D_{x,y}B(v,x,y)$, open at fresh $b$ and retain
$\Q_yB(s,b,y)$. Introduce the desired D formula at witness $b$.
From its base $B(t,b,0)$, reverse matrix transport gives $B(s,b,0)$,
QE gives $B(s,b,Sz)$, and forward transport gives $B(t,b,Sz)$.
DI$_c$ closes the desired formula and DE discharges $b$.
Each binder case uses only the strictly smaller matrix.
\end{proof}
This elementary fusion theorem assumes neither predecessor, numerical
equality tests, successor injectivity nor an induction-hypothesis rule.

